% ============================================================================
% Fichier  : partie1/P01_C03_controle_interne.tex
% Titre    : Le Contrôle Interne et ses Composantes
% Auteur   : MINKA MI NGUIDJOÏ Thierry Emmanuel
% Version  : 1.0 -- novembre 2026
% ============================================================================
% Base : ACSI C02 (faits LiZbiZ et SANAGA repris à l'identique : matrice des
% droits F1-F5, paires F1-F5, F2-F3, F3-F4, F2-F5 ; 4 utilisateurs sur 6 en
% conflit ; contrôles des stocks 1/10/1 ; SANAGA 214 contrôles dont 9
% informatiques, aucun sur les changements).
% Sources : COSO, Internal Control -- Integrated Framework (2013), cinq
% composantes et dix-sept principes paraphrasés ; règlement COBAC R-2016/04,
% articles 1 et 2 (texte relevé pour le cours GCI, à collationner) ;
% articles A0 et A1 pour la lecture µ, U, H.
% Corrigés : corriges/Corr_C03.tex.
% ============================================================================

\chapter{Le Contrôle Interne et ses Composantes}
\label{chap:controle-interne}

\epigraph{%
    « Un contrôle qui ne laisse pas de trace\\
    n'existe que dans la procédure qui le décrit. »%
}{Principe de l'auditeur}

\niveauChapitre{Fondamental}{%
    Maîtrise du Chapitre~\ref{chap:audit-auditologie} (audit, contrôle,
    trois lignes) et du Chapitre~\ref{chap:defaillances} (contrôle conçu
    mais inopérant)%
}

\begin{boxobjectifs}
\begin{itemize}
    \item Définir le contrôle interne, ses trois catégories d'objectifs et
          ses limites.
    \item Énoncer les cinq composantes et les dix-sept principes du
          COSO~2013, et illustrer chaque principe dans un système
          d'information.
    \item Classer un contrôle selon le moment, le mode d'exécution et la
          portée~; expliquer pourquoi les contrôles généraux informatiques
          conditionnent les contrôles d'application.
    \item Appliquer la séparation des tâches à une matrice des droits, et
          détecter les conflits par une requête SQL.
    \item Concevoir un dispositif à l'aide d'une matrice des risques et des
          contrôles complète.
    \item Évaluer la conception et l'efficacité d'un contrôle manuel et
          d'un contrôle automatisé, et coter le résultat \OUI{}, \NON{} ou
          \NEK{}.
\end{itemize}
\end{boxobjectifs}

\bigskip

% ============================================================================
\section*{Question directrice : où passaient les 186 millions~?}
% ============================================================================

L'écart de stock de LiZbiZ (Chapitre~\ref{chap:prologue}) ne s'explique
pas par un seul acte. Les premières investigations d'AUDIX CONSEIL montrent
que, dans l'application intégrée, un même magasinier peut enregistrer une
réception de marchandises \emph{et} ajuster le stock après inventaire. Le
responsable informatique dispose de tous les droits, y compris celui de
valider un paiement. Personne ne revoit ces droits.

Aucune règle n'interdisait ces cumuls~: LiZbiZ n'a jamais conçu son
contrôle interne, il s'est formé au fil des urgences. Les quatre affaires du
Chapitre~\ref{chap:defaillances} l'ont montré~: un contrôle peut exister
sur le papier et ne rien empêcher. Ce chapitre apprend à reconnaître un
dispositif conçu, à le concevoir soi-même et à l'évaluer.

\medskip
\textbf{Plan du chapitre.}
\begin{enumerate}
    \item La définition et les objectifs du contrôle interne
          (section~\ref{sec:c03-definition}).
    \item Le cadre COSO~2013, composante par composante
          (section~\ref{sec:c03-coso}).
    \item Le classement des contrôles et les contrôles généraux
          informatiques (sections~\ref{sec:c03-types}
          et~\ref{sec:c03-itgc}).
    \item La séparation des tâches, avec la détection automatique des
          conflits (section~\ref{sec:c03-sod}).
    \item La conception d'un dispositif (section~\ref{sec:c03-concevoir}).
    \item L'évaluation d'un contrôle et la lecture d'un domaine
          (section~\ref{sec:c03-evaluer}).
    \item Le contrôle interne des banques de la CEMAC
          (section~\ref{sec:c03-cobac}).
\end{enumerate}

\separateur

% ============================================================================
\section{Qu'est-ce que le Contrôle Interne~?}
\label{sec:c03-definition}
% ============================================================================

\subsection{Partir d'un exemple}

Chez LiZbiZ, un fournisseur livre 200~ordinateurs portables à l'entrepôt
de Douala. Trois choses doivent être vraies pour que la comptabilité
reflète la réalité~: les 200~ordinateurs ont bien été commandés~; ils ont
bien été reçus~; la facture payée correspond à ce qui a été reçu. Chacune
de ces trois affirmations peut être fausse, par erreur ou par fraude. Le
contrôle interne est l'ensemble des dispositions qui rendent ces
affirmations fiables~: un bon de commande validé, un bon de réception
signé par une autre personne que l'acheteur, un rapprochement de la
facture avec la commande et la réception avant le paiement.

Personne, chez LiZbiZ, n'appelle cela « contrôle interne ». C'est pourtant
ce que la direction attend lorsqu'elle signe les états financiers.

\subsection{La définition du COSO}

Le cadre de référence le plus répandu est celui du \sigle{COSO}
(\textit{Committee of Sponsoring Organizations of the Treadway Commission}),
dont la version actuelle date de 2013~\cite{coso2013ic}.

\begin{boxdef}
\termecle{Contrôle interne}\index{contrôle interne}~: processus mis en œuvre
par l'organe de gouvernance, la direction et l'ensemble du personnel, destiné
à fournir une assurance raisonnable quant à la réalisation d'objectifs liés
aux opérations, à l'information (\textit{reporting}) et à la conformité.
\end{boxdef}
{\small\itshape Définition reformulée d'après le COSO~2013~\cite{coso2013ic}.}

\medskip

Quatre idées sont contenues dans cette définition.
\begin{enumerate}
    \item \textbf{Un processus}, pas un document~: le contrôle interne vit
          dans les opérations quotidiennes. Un manuel de procédures que
          personne n'applique n'est pas un contrôle interne.
    \item \textbf{L'affaire de tous}~: il est porté par la direction et le
          personnel, pas seulement par le contrôleur interne.
    \item \textbf{Une assurance raisonnable}, jamais absolue
          (Chapitre~\ref{chap:audit-auditologie})~: le coût d'un contrôle
          ne doit pas dépasser le risque qu'il traite.
    \item \textbf{Des objectifs}~: on ne contrôle pas « en général », mais en
          vue d'objectifs précis. Sans objectif, il n'y a pas de risque, et
          sans risque, pas de contrôle à justifier.
\end{enumerate}

\subsection{Trois catégories d'objectifs}

\begin{tableaucours}{L{2.8cm}L{5.0cm}L{4.4cm}}{%
    \tableauHeader{Catégorie} &
    \tableauHeader{Question} &
    \tableauHeader{Exemple chez LiZbiZ}
}
Opérations & Les activités sont-elles efficaces et efficientes, les actifs
protégés~? & Le stock enregistré correspond au stock réel \\
Information & L'information financière et non financière est-elle
fiable~? & Les états financiers reflètent la valeur du stock \\
Conformité & Les lois et règlements sont-ils respectés~? & Les données
personnelles des clients sont protégées (loi n°~2024/017) \\
\end{tableaucours}
\vspace{-0.8em}
\begin{center}{\small\itshape Tableau~3.1. Les objectifs du contrôle
interne}\end{center}

\subsection{Les assertions}

Un objectif d'information se décompose en \termecle{assertions}%
\index{assertion}, c'est-à-dire en affirmations implicites de la direction
sur les données. L'auditeur les reprend telles quelles~: chaque contrôle
protège une ou plusieurs assertions.

\begin{tableaucours}{L{2.8cm}L{5.4cm}L{4.0cm}}{%
    \tableauHeader{Assertion} &
    \tableauHeader{Ce qu'elle affirme} &
    \tableauHeader{Contrôle typique}
}
Réalité & Les opérations enregistrées ont eu lieu & Bon de réception signé
par un tiers \\
Exhaustivité & Toutes les opérations sont enregistrées & Séquence continue
des numéros de bons \\
Exactitude & Les montants et quantités sont justes & Rapprochement
commande, réception, facture \\
Séparation des exercices & L'opération est rattachée à la bonne période &
Contrôle des réceptions de fin d'année \\
Droits et obligations & L'entité détient l'actif ou doit la dette &
Contrat, titre de propriété \\
Présentation & L'information est classée et décrite correctement & Revue
des états par la direction financière \\
\end{tableaucours}
\vspace{-0.8em}
\begin{center}{\small\itshape Tableau~3.2. Les assertions et les
contrôles qui les protègent}\end{center}

\subsection{Ce que le contrôle interne n'est pas}

\begin{itemize}
    \item \textbf{Ce n'est pas l'audit interne.} Le contrôle interne est
          exécuté par les opérationnels (première ligne) et supervisé par
          les fonctions de risque et de conformité (deuxième ligne)~;
          l'audit interne (troisième ligne) l'évalue
          (Chapitre~\ref{chap:audit-auditologie}).
    \item \textbf{Ce n'est pas la sécurité informatique}, même si une
          partie des contrôles sont des mesures de sécurité. Le contrôle
          interne couvre aussi la commande, la comptabilité, la paie.
    \item \textbf{Ce n'est pas une garantie.} Il réduit la probabilité et
          l'effet des anomalies~; il ne les supprime pas
          (section~\ref{sec:c03-limites}).
\end{itemize}

% ============================================================================
\section{Le Cadre COSO 2013}
\label{sec:c03-coso}
% ============================================================================

\subsection{Cinq composantes, dix-sept principes}

Le COSO organise le contrôle interne en \textbf{cinq composantes}, détaillées
en \textbf{dix-sept principes}. Un dispositif n'est efficace que si les cinq
composantes sont présentes, fonctionnent et agissent ensemble~: une
activité de contrôle excellente ne compense pas une direction qui passe
outre.

\begin{tableaucours}{L{3.2cm}cL{8.0cm}}{%
    \tableauHeader{Composante} &
    \tableauHeader{N°} &
    \tableauHeader{Principe (reformulé)}
}
Environnement de contrôle & 1 & Afficher un engagement envers l'intégrité et
l'éthique \\
 & 2 & Faire superviser le contrôle interne par un organe de gouvernance
indépendant de la direction \\
 & 3 & Définir les structures, les rattachements, les pouvoirs et les
responsabilités \\
 & 4 & Attirer, développer et retenir des personnes compétentes \\
 & 5 & Faire rendre compte chacun de ses responsabilités de contrôle \\
Évaluation des risques & 6 & Fixer des objectifs suffisamment clairs pour
identifier les risques \\
 & 7 & Identifier et analyser les risques qui menacent ces objectifs \\
 & 8 & Prendre en compte le risque de fraude \\
 & 9 & Identifier et analyser les changements significatifs \\
Activités de contrôle & 10 & Choisir et mettre en place des activités de
contrôle qui réduisent les risques \\
 & 11 & Choisir et mettre en place des contrôles généraux sur les
technologies \\
 & 12 & Déployer les contrôles par des politiques et des procédures \\
Information et communication & 13 & Produire et utiliser une information
pertinente et de qualité \\
 & 14 & Communiquer en interne sur le contrôle interne \\
 & 15 & Communiquer avec les tiers \\
Pilotage & 16 & Conduire des évaluations continues ou ponctuelles \\
 & 17 & Évaluer les déficiences et les communiquer à temps \\
\end{tableaucours}
\vspace{-0.8em}
\begin{center}{\small\itshape Tableau~3.3. Les composantes et les principes
du COSO~2013 (notre reformulation)}\end{center}

\subsection{Les principes vus depuis le système d'information}

Le COSO est écrit pour toute l'organisation. L'auditeur des SI le lit avec
une question~: qu'est-ce que chaque principe exige du système
d'information~? Le Tableau~3.4 donne, pour chaque principe, une preuve que
l'auditeur peut demander.

\begin{tableaucours}{cL{5.6cm}L{6.0cm}}{%
    \tableauHeader{N°} &
    \tableauHeader{Ce que le principe exige du SI} &
    \tableauHeader{Preuve que l'auditeur peut demander}
}
1 & Une charte d'usage des ressources informatiques signée & Charte,
accusés de réception des salariés \\
2 & Un comité qui examine les risques informatiques & Procès-verbaux du
comité d'audit ou des risques \\
3 & Des rôles écrits pour la DSI, le RSSI, les propriétaires
d'application & Organigramme, fiches de poste, délégations \\
4 & Des compétences informatiques maintenues & Plan de formation,
certifications \\
5 & Des propriétaires d'application qui valident les droits & Revues des
droits signées \\
6 & Des objectifs de disponibilité et de sécurité chiffrés & Engagements
de service, politique de sécurité \\
7 & Une analyse des risques informatiques & Registre des risques
(Chapitre~\ref{chap:gestion-risques}) \\
8 & Une analyse des fraudes possibles par le système & Scénarios de fraude
(Chapitre~\ref{chap:fraude}) \\
9 & L'analyse des changements majeurs (nouveau système, migration) &
Dossier de décision de mise en production \\
10 & Des contrôles d'application dans les processus & Paramétrage,
contrôles bloquants \\
11 & Des contrôles généraux informatiques & Gestion des accès, des
changements, de l'exploitation (section~\ref{sec:c03-itgc}) \\
12 & Des procédures informatiques écrites et appliquées & Procédures,
preuves d'exécution \\
13 & Des données de pilotage fiables & Règles de qualité des données \\
14 & Une communication interne sur la sécurité & Campagnes, intranet \\
15 & Des informations fiables échangées avec les tiers & Contrats,
conventions d'échange \\
16 & Une surveillance continue des contrôles & Tableaux de bord, alertes
(Chapitre~\ref{chap:audit-continu}) \\
17 & Un suivi des déficiences informatiques & Plan d'action, suivi des
recommandations \\
\end{tableaucours}
\vspace{-0.8em}
\begin{center}{\small\itshape Tableau~3.4. Les dix-sept principes et le
système d'information}\end{center}

\begin{boxinfo}
\textbf{Le principe~11 concerne directement l'auditeur des SI.} Le COSO
demande de choisir et de mettre en place des contrôles généraux sur les
technologies~: accès, changements, exploitation. C'est le cœur des
contrôles généraux informatiques (section~\ref{sec:c03-itgc}).
\end{boxinfo}

\subsection{Présent, fonctionnel, intégré}

Le COSO juge un dispositif efficace lorsque chaque composante et chaque
principe sont \textbf{présents} (conçus), \textbf{fonctionnels}
(exécutés) et que les cinq composantes \textbf{agissent ensemble}. Une
\termecle{déficience majeure}\index{déficience majeure} est une déficience
qui réduit significativement la probabilité d'atteindre les objectifs~;
elle interdit de conclure à l'efficacité du dispositif. Le
Chapitre~\ref{chap:referentiels} appliquera cette évaluation aux dix-sept
principes de la BANQUE SANAGA.

\subsection{Le diagnostic de LiZbiZ}

\begin{boxscenario}{Opération LIZBIZ. Les cinq composantes}
\textbf{Constatations} (entretiens et documents remis à AUDIX CONSEIL)~:
\begin{itemize}
    \item \emph{Environnement de contrôle}~: pas de comité d'audit~; les
          pouvoirs de validation ne sont écrits nulle part.
    \item \emph{Évaluation des risques}~: aucune cartographie des risques~;
          le risque de fraude n'a jamais été examiné.
    \item \emph{Activités de contrôle}~: des contrôles existent dans les
          agences, mais varient d'un site à l'autre.
    \item \emph{Information et communication}~: les écarts d'inventaire ne
          remontent à la direction qu'une fois par an.
    \item \emph{Pilotage}~: aucune évaluation du dispositif depuis la
          création de l'entreprise.
\end{itemize}

\textbf{Analyse} (confiance moyenne). Les activités de contrôle ne peuvent
pas compenser l'absence des quatre autres composantes. L'écart de stock est
le symptôme d'un dispositif qui n'a jamais été conçu.
\end{boxscenario}

% ============================================================================
\section{Classer les Contrôles}
\label{sec:c03-types}
% ============================================================================

\subsection{Selon le moment}

\begin{boxdef}
\termecle{Contrôle préventif}\index{contrôle préventif}~: empêche
l'anomalie de se produire (une validation obligatoire avant paiement).

\medskip
\termecle{Contrôle détectif}\index{contrôle détectif}~: repère l'anomalie
après coup (le rapprochement mensuel du stock et des inventaires).

\medskip
\termecle{Contrôle correctif}\index{contrôle correctif}~: corrige l'anomalie
et ses effets (la restauration d'une sauvegarde après une suppression).
\end{boxdef}

Un dispositif solide combine les trois. Le préventif est préférable, car il
évite le dommage~; mais il peut être contourné, et seul le détectif révèle
alors le contournement.

\subsection{Selon le mode d'exécution}

Un contrôle peut être \textbf{manuel} (une signature sur un bon de sortie),
\textbf{automatisé} (l'application refuse un ajustement de stock non
validé) ou \textbf{semi-automatisé}, c'est-à-dire dépendant de
l'informatique (l'application produit une liste d'anomalies qu'une personne
examine).

\begin{tableaucours}{L{2.8cm}L{4.6cm}L{4.8cm}}{%
    \tableauHeader{Mode} &
    \tableauHeader{Forces} &
    \tableauHeader{Faiblesses}
}
Manuel & Jugement humain, adaptable & Variable, oublié en période de
charge, preuve souvent absente \\
Automatisé & Constant, exhaustif, rapide & Ne vaut que si le paramétrage
est protégé~; une erreur se répète à chaque opération \\
Semi-automatisé & Combine exhaustivité et jugement & Dépend à la fois de
l'exactitude de la liste et de la diligence de la personne \\
\end{tableaucours}
\vspace{-0.8em}
\begin{center}{\small\itshape Tableau~3.5. Les modes d'exécution d'un
contrôle}\end{center}

\subsection{Selon la portée}

\begin{boxdef}
\termecle{Contrôles généraux informatiques}\index{contrôles généraux
informatiques} (\textit{IT general controls}, \sigle{ITGC})~: contrôles qui
s'appliquent à l'environnement dans lequel les applications sont
développées, exploitées et maintenues~: gestion des accès, gestion des
changements, exploitation (sauvegardes, traitements, incidents).

\medskip
\termecle{Contrôles d'application}\index{contrôles d'application}~:
contrôles intégrés à une application, qui agissent sur le traitement des
opérations~: contrôles des entrées (validation des données), des traitements
(exactitude des calculs) et des sorties (complétude des états).
\end{boxdef}

\begin{boxalert}
\textbf{Violation typique.} L'auditeur constate que l'application refuse tout
paiement d'un montant supérieur au plafond autorisé, et conclut que le
contrôle est efficace. Or le responsable informatique peut modifier le
plafond directement dans la base, sans trace.

\smallskip
\textbf{Règle~:} un contrôle d'application ne vaut que si les contrôles
généraux qui le protègent (accès, changements) sont eux-mêmes efficaces.
L'auditeur les teste en premier.
\end{boxalert}

% ============================================================================
\section{Les Contrôles Généraux Informatiques}
\label{sec:c03-itgc}
% ============================================================================

\subsection{Quatre familles}

Les contrôles généraux se regroupent en quatre familles, qui reviennent
dans tous les référentiels (COBIT, ISO/IEC~27002, guide du CHAI) sous des
noms différents.

\begin{tableaucours}{L{3.0cm}L{4.8cm}L{4.4cm}}{%
    \tableauHeader{Famille} &
    \tableauHeader{Objectif} &
    \tableauHeader{Contrôles typiques}
}
Accès & Seules les personnes autorisées accèdent aux programmes et aux
données, avec les seuls droits nécessaires & Création et retrait des
comptes, revue des droits, comptes à privilèges, authentification \\
Changements & Aucune modification de programme ou de paramètre n'est mise
en production sans être demandée, testée et approuvée & Demande, tests,
approbation, séparation des environnements, journal des mises en
production \\
Exploitation & Les traitements s'exécutent comme prévu et les données
sont récupérables & Ordonnancement, sauvegardes, restaurations, gestion
des incidents \\
Développement et acquisition & Les nouveaux systèmes satisfont aux
exigences de contrôle & Exigences, recette, migration des données
(Chapitre~\ref{chap:audit-projets}) \\
\end{tableaucours}
\vspace{-0.8em}
\begin{center}{\small\itshape Tableau~3.6. Les quatre familles de
contrôles généraux informatiques}\end{center}

\subsection{Pourquoi l'ordre compte}

Supposons que l'application de LiZbiZ bloque les paiements supérieurs à
5~millions de FCFA sans seconde signature. Trois questions précèdent la
conclusion que ce contrôle est efficace~:
\begin{enumerate}
    \item \textbf{Accès}~: qui peut modifier le seuil de 5~millions~? Si
          c'est un administrateur sans revue, le seuil peut être relevé la
          veille d'un paiement et rabaissé le lendemain.
    \item \textbf{Changements}~: le programme qui applique le seuil a-t-il
          été modifié pendant la période~? Une modification non testée peut
          avoir désactivé la règle.
    \item \textbf{Exploitation}~: les paiements passent-ils tous par
          l'application, ou certains sont-ils saisis directement dans la
          banque en ligne~?
\end{enumerate}
Si l'une des trois réponses est défavorable, le contrôle d'application ne
peut pas être coté \OUI{}, même s'il fonctionne sur les opérations
examinées.

\subsection{Tester un contrôle automatisé}

Un contrôle automatisé se teste autrement qu'un contrôle manuel. Comme il
s'exécute à l'identique à chaque opération, un seul cas bien choisi suffit
à établir qu'il fonctionne, \emph{à condition} que les contrôles généraux
garantissent qu'il n'a pas changé pendant la période.

\begin{boxprocedure}
\textbf{Tester le seuil de double signature (contrôle automatisé)}
\begin{enumerate}
    \item \textbf{Lire le paramétrage}~: capture d'écran ou extraction de
          la table des paramètres, horodatée, montrant le seuil de
          5~000~000~FCFA.
    \item \textbf{Vérifier l'historique}~: journal des modifications du
          paramètre sur la période (aucune modification, ou modifications
          autorisées et tracées).
    \item \textbf{Test positif}~: dans l'environnement de test, saisir un
          paiement de 5~000~001~FCFA avec une seule signature~; constater le
          refus.
    \item \textbf{Test négatif}~: saisir un paiement de 4~999~999~FCFA~;
          constater qu'il passe avec une seule signature.
    \item \textbf{Vérifier l'exhaustivité}~: rapprocher le total des
          paiements de l'application et le total des relevés bancaires, pour
          s'assurer qu'aucun paiement ne contourne l'application.
    \item \textbf{Conserver les preuves} (captures, extraction, rapprochement)
          dans les papiers de travail.
\end{enumerate}
\end{boxprocedure}

\begin{boxinfo}
\textbf{Ce que le test ne dit pas.} Le seuil empêche un paiement unique
supérieur à 5~millions~; il n'empêche pas vingt-trois paiements de
4,9~millions au même fournisseur. Le Chapitre~\ref{chap:fraude} montrera
que c'est exactement ce qui s'est passé chez LiZbiZ.
\end{boxinfo}

% ============================================================================
\section{Séparer les Tâches}
\label{sec:c03-sod}
% ============================================================================

\subsection{Le principe}

\begin{boxdef}
\termecle{Séparation des tâches}\index{séparation des tâches}
(\textit{segregation of duties}, \sigle{SoD})~: répartition, entre des
personnes différentes, des fonctions dont le cumul permettrait de commettre
et de dissimuler une erreur ou une fraude. On sépare en principe quatre
fonctions~: \emph{autoriser} une opération, l'\emph{enregistrer},
\emph{garder} l'actif concerné, et \emph{contrôler}.

\medskip
\termecle{Règle des quatre yeux}\index{règle des quatre yeux}~: toute
opération sensible est réalisée par une personne et validée par une autre.
\end{boxdef}

\subsection{La matrice des droits de LiZbiZ}

\begin{boxscenario}{Opération LIZBIZ. Qui peut faire quoi~?}
AUDIX CONSEIL extrait de l'application intégrée les droits de six
utilisateurs sur cinq fonctions sensibles~: F1 créer ou modifier un
fournisseur~; F2 saisir une commande d'achat~; F3 enregistrer une réception
en entrepôt~; F4 ajuster le stock après inventaire~; F5 valider un paiement
fournisseur.

\medskip
\begin{tableaucours}{L{4.6cm}ccccc}{%
    \tableauHeader{Utilisateur (fonction)} &
    \tableauHeader{F1} & \tableauHeader{F2} & \tableauHeader{F3} &
    \tableauHeader{F4} & \tableauHeader{F5}
}
M.~KOUNOU (resp. informatique) & × & × & × & × & × \\
M.~MBIDA (acheteur) & × & × & & & \\
M.~FOUDA (magasinier, Douala) & & & × & × & \\
Mme~OWONA (comptable) & × & & & & × \\
M.~NANGA (chef d'agence, Garoua) & & × & × & & \\
Mme~MBALLA (DAF) & & & & & × \\
\end{tableaucours}

\medskip
\textbf{Paires incompatibles retenues.} F1 et F5 (créer un fournisseur
fictif, puis le payer)~; F2 et F3 (commander, puis attester soi-même une
réception qui n'a pas eu lieu)~; F3 et F4 (enregistrer une réception, puis
« corriger » le stock pour masquer un manquant)~; F2 et F5 (commander, puis
payer sa propre commande).

\medskip
\textbf{Constatation.} Quatre utilisateurs sur six cumulent au moins une
paire incompatible~: M.~KOUNOU (toutes), M.~FOUDA (F3 et F4), Mme~OWONA
(F1 et F5) et M.~NANGA (F2 et F3).

\medskip
\textbf{Analyse} (confiance élevée). Le cumul F3 et F4 du magasinier de
Douala suffit à expliquer qu'un manquant puisse disparaître des écritures.
Le cumul total du responsable informatique est le plus grave~: il cumule les
droits de gestion et le pouvoir technique de modifier les droits eux-mêmes.
\end{boxscenario}

\subsection{Détecter les conflits par une requête}

Sur six utilisateurs, la matrice se lit à l'œil. Sur les 150~salariés de
LiZbiZ et les quelque 40~fonctions de l'application, ce n'est plus
possible. L'auditeur demande trois extractions (utilisateurs, droits,
fonctions) et formalise les paires incompatibles dans une table.

\begin{boxoutils}{Requête SQL de détection des conflits}
\begin{lstlisting}[language=SQL]
-- droits(utilisateur, fonction) : extraction de l'application
-- paires(f_a, f_b, motif)       : paires incompatibles, écrites
--                                 par l'auditeur et validées
SELECT d1.utilisateur, p.f_a, p.f_b, p.motif
FROM   droits d1
JOIN   droits d2 ON d1.utilisateur = d2.utilisateur
JOIN   paires p  ON p.f_a = d1.fonction AND p.f_b = d2.fonction
ORDER  BY d1.utilisateur, p.f_a;
\end{lstlisting}
Sur la matrice de LiZbiZ, avec les quatre paires retenues, la requête rend
sept lignes~: quatre pour M.~KOUNOU, une pour M.~FOUDA, une pour
Mme~OWONA, une pour M.~NANGA. M.~MBIDA n'apparaît pas~: son cumul F1 et F2
(créer un fournisseur et lui commander) n'est pas une paire retenue par
l'entreprise, ce que l'auditeur peut discuter.
\end{boxoutils}

\begin{boxalert}
\textbf{Violation typique.} Exécuter la requête sur la seule table des
droits, sans vérifier que l'extraction est complète (comptes désactivés,
comptes techniques, droits accordés par un rôle plutôt qu'à la personne).

\smallskip
\textbf{Règle~:} une extraction se rapproche d'un total indépendant (le
nombre de comptes affiché par l'outil d'administration) et se fait en
présence de l'auditeur ou par une requête qu'il fournit. Les droits
hérités d'un rôle sont « dépliés » avant l'analyse.
\end{boxalert}

\subsection{Quand on ne peut pas séparer}

Dans une petite agence de trois personnes, une séparation complète est
impossible. On recourt alors à des \termecle{contrôles
compensatoires}\index{contrôle compensatoire}~: revue a posteriori par une
personne extérieure à l'agence, rapprochement mensuel par le siège, journal
des ajustements examiné par le contrôleur interne. Le contrôle compensatoire
doit être écrit, exécuté et tracé, comme tout autre contrôle.

\begin{boxmetier}{Contrôleur interne}
\textbf{Mettre en place le contrôle des stocks chez LiZbiZ}
\begin{enumerate}
    \item Retirer à M.~FOUDA le droit F4~; confier les ajustements de stock
          à une personne qui ne gère pas les réceptions.
    \item Rendre l'ajustement soumis à validation (quatre yeux) au-delà d'un
          seuil, par exemple 500~000~FCFA.
    \item Retirer au responsable informatique tous les droits de gestion~;
          tracer ses accès d'administration.
    \item Instaurer un inventaire tournant mensuel dans chaque entrepôt, et
          un rapprochement par le siège.
    \item Consigner le tout dans la matrice des risques et des contrôles, et
          prévoir une revue trimestrielle des droits.
\end{enumerate}
\end{boxmetier}

% ============================================================================
\section{Concevoir un Dispositif}
\label{sec:c03-concevoir}
% ============================================================================

\subsection{Six étapes}

Mettre en place un contrôle interne, c'est relier chaque risque à un ou
plusieurs contrôles, et organiser la preuve de leur exécution.

\begin{enumerate}
    \item \textbf{Fixer les objectifs} du processus (par exemple~: le stock
          enregistré est exact et complet).
    \item \textbf{Identifier les risques} qui menacent ces objectifs
          (Chapitre~\ref{chap:risque}).
    \item \textbf{Décrire le processus} et ses points de risque.
    \item \textbf{Choisir les contrôles}~: préventifs d'abord, détectifs en
          complément, automatisés lorsque c'est possible.
    \item \textbf{Documenter} chaque contrôle dans une matrice des risques
          et des contrôles.
    \item \textbf{Tester} puis \textbf{piloter} le dispositif.
\end{enumerate}

\begin{boxprocedure}
\textbf{La matrice des risques et des contrôles.} Pour que le dispositif soit
auditable, chaque ligne de la matrice indique~:
\begin{itemize}
    \item le processus, l'objectif, l'assertion et le risque traité~;
    \item le contrôle, sa nature (préventif ou détectif, manuel ou
          automatisé) et sa fréquence~;
    \item le responsable de l'exécution~;
    \item \textbf{la preuve} que laisse chaque exécution (signature,
          journal, état signé)~;
    \item la date de la dernière évaluation et son résultat.
\end{itemize}
Un contrôle sans preuve d'exécution ne peut pas être évalué~: il doit être
reconçu.
\end{boxprocedure}

\subsection{Une matrice pour les stocks de LiZbiZ}

Le Tableau~3.7 donne quatre lignes de la matrice proposée par AUDIX CONSEIL
pour le processus de gestion des stocks.

\begin{tableaucours}{L{2.6cm}L{3.2cm}L{2.6cm}L{3.8cm}}{%
    \tableauHeader{Risque} &
    \tableauHeader{Contrôle} &
    \tableauHeader{Nature, fréquence} &
    \tableauHeader{Responsable et preuve}
}
Réception fictive (réalité) & Réception saisie par un magasinier distinct
de l'acheteur & Préventif, automatisé, permanent & Paramétrage des
droits~; revue trimestrielle signée par le DAF \\
Manquant masqué par un ajustement (exactitude) & Ajustement au-delà de
500~000~FCFA validé par le contrôleur de gestion & Préventif, automatisé,
à chaque ajustement & Journal des validations dans l'application \\
Ajustements fractionnés (exactitude) & Revue mensuelle des ajustements
par magasinier et par jour & Détectif, semi-automatisé, mensuel & État
d'anomalies annoté et signé par la contrôleuse interne \\
Stock non compté (exhaustivité) & Inventaire tournant mensuel, 1/12 des
références par mois & Détectif, manuel, mensuel & Feuilles de comptage
signées par deux personnes, rapprochement par le siège \\
\end{tableaucours}
\vspace{-0.8em}
\begin{center}{\small\itshape Tableau~3.7. Extrait de la matrice des
risques et des contrôles des stocks}\end{center}

\begin{boxinfo}
\textbf{Pourquoi la troisième ligne.} Le seuil de 500~000~FCFA de la
deuxième ligne invite au fractionnement~: un ajustement de 1,2~million se
découpe en trois de 400~000. La revue mensuelle par magasinier et par jour
est le contrôle détectif qui compense cette faiblesse connue du contrôle
préventif. Concevoir un dispositif, c'est aussi anticiper le contournement
de ses propres contrôles.
\end{boxinfo}

% ============================================================================
\section{Évaluer un Contrôle}
\label{sec:c03-evaluer}
% ============================================================================

\subsection{Conception et efficacité}

L'auditeur évalue un contrôle en deux temps.

\begin{tableaucours}{L{3.2cm}L{4.6cm}L{4.4cm}}{%
    \tableauHeader{Évaluation} &
    \tableauHeader{Question} &
    \tableauHeader{Comment}
}
Conception & Le contrôle, s'il fonctionne, réduit-il le risque~? &
Entretien, lecture de la procédure, test de cheminement sur un cas \\
Efficacité opérationnelle & Le contrôle a-t-il fonctionné comme prévu
pendant toute la période~? & Test d'un échantillon d'exécutions et de leurs
preuves \\
\end{tableaucours}
\vspace{-0.8em}
\begin{center}{\small\itshape Tableau~3.8. Deux niveaux d'évaluation d'un
contrôle}\end{center}

Un contrôle mal conçu est inefficace, même s'il est exécuté à la lettre. Un
contrôle bien conçu mais non exécuté ne protège de rien.

\begin{boxdef}
\termecle{Test de cheminement}\index{test de cheminement}
(\textit{walkthrough})~: suivi d'une opération réelle, du début à la fin du
processus, en examinant à chaque étape le document produit et le contrôle
exercé. Il établit que l'auditeur a compris le processus et que le contrôle
existe tel qu'il est décrit~; il ne suffit pas à établir qu'il a fonctionné
toute la période.
\end{boxdef}

\subsection{Combien d'exécutions tester~?}

L'efficacité d'un contrôle manuel se teste sur un échantillon dont la taille
dépend de la fréquence du contrôle. Le Tableau~3.9 donne les ordres de
grandeur indicatifs utilisés dans ce cours~; le Chapitre~\ref{chap:preuve}
en donnera le fondement statistique.

\begin{tableaucours}{L{3.4cm}L{3.4cm}L{5.4cm}}{%
    \tableauHeader{Fréquence du contrôle} &
    \tableauHeader{Exécutions par an} &
    \tableauHeader{Échantillon indicatif}
}
Annuelle & 1 & 1 \\
Trimestrielle & 4 & 2 \\
Mensuelle & 12 & 2 à 5 \\
Hebdomadaire & 52 & 5 à 15 \\
Quotidienne & environ 250 & 20 à 40 \\
Plusieurs fois par jour & plus de 250 & 25 à 60 \\
Automatisée & & 1 cas positif et 1 négatif, si les contrôles généraux
sont efficaces \\
\end{tableaucours}
\vspace{-0.8em}
\begin{center}{\small\itshape Tableau~3.9. Tailles d'échantillon
indicatives (convention du cours)}\end{center}

\subsection{\NON{} ou \NEK{}~?}

La cotation d'un contrôle demande une distinction que l'audit binaire ne
fait pas.

\begin{itemize}
    \item \textbf{\NON{}}~: l'auditeur établit que le contrôle n'est pas
          conforme au critère. Par exemple, le critère exige une validation
          tracée de chaque ajustement, et l'extraction montre des
          ajustements sans validation.
    \item \textbf{\NEK{}}~: l'information disponible ne permet pas de
          conclure, malgré des diligences normales. Par exemple, le journal
          des ajustements n'a pas été remis, ou deux sources se contredisent
          sans qu'on puisse les départager.
\end{itemize}

\begin{boxalert}
\textbf{Violation typique.} Coter \NEK{} un contrôle dont le critère exige une
trace, alors que l'auditeur a vérifié qu'aucune trace n'existe.

\smallskip
\textbf{Règle~:} l'absence d'une trace exigée par le critère est un \NON{}
établi. \NEK{} est réservé au cas où l'auditeur ne peut pas savoir~; il
s'accompagne toujours de la description des diligences effectuées.
\end{boxalert}

\begin{boxobjet}{Ajustements de stock validés}
\textbf{Objectif.} Les ajustements de stock reflètent des écarts
d'inventaire réels.

\textbf{Critère.} Tout ajustement supérieur à 500~000~FCFA est validé dans
l'application par une personne autre que son auteur (matrice des risques et
des contrôles, ligne~2).

\textbf{Question.} Chaque ajustement de plus de 500~000~FCFA de la période
porte-t-il une validation par une autre personne~?

\textbf{Risque.} Un manquant (vol, casse, erreur) est effacé des écritures
sans examen.

\textbf{Conséquence.} Le stock et le résultat sont surévalués~; une fraude
peut se poursuivre sans être vue.

\textbf{Test.} Extraction de tous les ajustements de janvier à avril~2026~;
filtre au-delà du seuil~; comparaison de l'auteur et du valideur.

\textbf{Preuve et trace.} Extraction horodatée et son empreinte SHA-256~;
requête jointe aux papiers de travail.
\end{boxobjet}

\subsection{Lire un domaine}

\begin{boxlecture}{Contrôles des stocks de LiZbiZ}
Douze contrôles de la gestion des stocks sont évalués sur la période de
janvier à avril~2026~: 1~\OUI{}, 10~\NON{}, 1~\NEK{}.

\medskip
\textbf{Constatation.} $\mu = 1/12 \approx 0{,}08$~; $U = 1/12 \approx 0{,}08$~;
$H \approx 0{,}82$~bit.

\medskip
\textbf{Analyse} (confiance élevée). La conformité est très faible, mais
l'incertitude l'est aussi~: l'organisation sait que ses contrôles de stock
ne fonctionnent pas. C'est une \emph{non-conformité connue}, qui appelle des
correctifs directs, à l'inverse du domaine~D du prologue ($H \approx 1{,}24$),
où il fallait d'abord comprendre.

\medskip
\hfill\lecture{0{,}08}{0{,}08}{0{,}82}
\end{boxlecture}

\begin{boxinfo}
\textbf{Le calcul de $H$, posé.} Avec $p_{\OUI} = 1/12$,
$p_{\NON} = 10/12$, $p_{\NEK} = 1/12$~:
\[
H = -\tfrac{1}{12}\log_2\tfrac{1}{12} - \tfrac{10}{12}\log_2\tfrac{10}{12}
    - \tfrac{1}{12}\log_2\tfrac{1}{12}
  \approx 0{,}299 + 0{,}219 + 0{,}299 \approx 0{,}82~\text{bit.}
\]
Le maximum, pour trois réponses, est $\log_2 3 \approx 1{,}58$~bit.
\end{boxinfo}

\subsection{Les limites du contrôle interne}
\label{sec:c03-limites}

Même bien conçu, un dispositif ne donne qu'une assurance raisonnable. Il
peut être mis en échec par la collusion entre plusieurs personnes, par un
dirigeant qui passe outre, par une erreur de jugement, ou parce que le coût
d'un contrôle dépasserait le risque qu'il traite. L'auditeur en tient compte
en recherchant, au-delà des contrôles, les indices de leur contournement
(Chapitre~\ref{chap:fraude}).

% ============================================================================
\section{Le Contrôle Interne dans les Banques}
\label{sec:c03-cobac}
% ============================================================================

Pour les établissements de crédit de la CEMAC, le contrôle interne n'est
pas une option~: c'est une obligation réglementaire~\cite{cobac2016r04}.

\begin{boxjuris}
\textbf{Règlement COBAC R-2016/04, articles~1 et~2} (texte relevé sur une
base réglementaire en ligne, à collationner avec le texte officiel)~:

\medskip
Article~1~: \textit{« Le présent règlement fixe les exigences minimales
relatives au système de contrôle interne que les établissements assujettis
doivent mettre en œuvre. »}

\medskip
Article~2, définition du risque opérationnel~: \textit{« risque de pertes
résultant de l'inadéquation ou de la défaillance des processus internes,
des personnes et des systèmes, ou d'évènements externes. »}

\medskip
\textbf{Analyse.} La défaillance d'un système informatique est, par
définition, un risque opérationnel que le dispositif de contrôle interne de
la banque doit maîtriser. L'auditeur des SI d'une banque teste donc les
contrôles informatiques au titre de ce règlement, et pas seulement au titre
des bonnes pratiques.
\end{boxjuris}

\begin{boxscenario}{Opération SANAGA. Le contrôle interne après la panne}
\textbf{Faits.} Après la panne de trois jours de mars~2026, la COBAC demande
à la BANQUE SANAGA de démontrer que son dispositif de contrôle interne
couvre les risques liés aux systèmes d'information.

\medskip
\textbf{Constatation} (audit interne, Mme~EKAMBI). La matrice des risques
et des contrôles de la banque compte 214~contrôles, dont 9~seulement portent
sur l'informatique, et aucun sur la gestion des changements.

\medskip
\textbf{Analyse} (confiance élevée). Le risque qui s'est réalisé, un
changement non maîtrisé du système central, n'était couvert par aucun
contrôle~: le défaut est de conception, pas d'exécution. Le
Chapitre~\ref{chap:audit-continuite} reviendra sur l'audit de la continuité
de la banque.
\end{boxscenario}

\begin{boxmetier}{Gestionnaire de risques}
\textbf{Relier le registre des risques et la matrice des contrôles.} Le
gestionnaire de risques de la BANQUE SANAGA vérifie que chaque risque
informatique de son registre renvoie à au moins une ligne de la matrice des
contrôles, et inversement. Un risque sans contrôle est un trou de
conception~; un contrôle sans risque est un coût sans justification.
\end{boxmetier}

\begin{boxmetier}{Auditeur SI}
\textbf{Commencer par les contrôles généraux.} Avant de tester un contrôle
d'application de la banque, l'auditeur obtient la liste des mises en
production de la période et la revue des comptes à privilèges du système
central. Sans elles, aucun contrôle automatisé ne peut être coté \OUI{}.
\end{boxmetier}

\begin{boxinfo}
\textbf{Ce que dit la recherche.} L'article A0~\cite{minka2026auditologie}
fait de l'observabilité partielle le premier axiome de l'audit~: l'auditeur
ne voit jamais tout le dispositif, il en voit des traces. La règle « une
trace exigée et absente donne \NON{} » en est la traduction
opérationnelle~: l'absence d'une trace qui devait exister est une
information, pas une ignorance.
\end{boxinfo}

\separateur

\begin{boxresume}
\textbf{Ce qu'il faut retenir de ce chapitre~:}
\begin{itemize}
    \item \textbf{Le contrôle interne est un processus}, porté par tous,
          qui donne une assurance raisonnable sur les opérations,
          l'information et la conformité~; il protège des assertions.
    \item \textbf{Cinq composantes, dix-sept principes} (COSO~2013)~: les
          activités de contrôle ne suffisent pas sans environnement,
          évaluation des risques, information et pilotage. Chaque principe
          a sa traduction dans le SI.
    \item \textbf{Un contrôle se classe} selon le moment, le mode et la
          portée~; les contrôles généraux (accès, changements,
          exploitation, développement) conditionnent les contrôles
          d'application.
    \item \textbf{Un contrôle automatisé} se teste sur un cas positif et un
          cas négatif, si les contrôles généraux sont efficaces.
    \item \textbf{La séparation des tâches} interdit le cumul des fonctions
          qui permettraient de commettre et de dissimuler~; les conflits se
          détectent par une requête sur des extractions complètes~; à
          défaut de séparation, des contrôles compensatoires tracés.
    \item \textbf{Un dispositif se conçoit} dans une matrice des risques et
          des contrôles, où chaque contrôle laisse une preuve.
    \item \textbf{Évaluer un contrôle}, c'est juger sa conception puis son
          efficacité~; l'absence d'une trace exigée est un \NON{}, pas un
          \NEK{}.
\end{itemize}
\end{boxresume}

\begin{boxinfo}
\textbf{Points de contrôle de l'élève.} Je sais~: citer les cinq
composantes du COSO~; classer un contrôle selon trois axes~; nommer les
quatre familles de contrôles généraux~; écrire la requête de détection des
conflits~; remplir une ligne de matrice des risques et des contrôles~;
distinguer un \NON{} d'un \NEK{}.
\end{boxinfo}

\bigskip

\begin{boxexo}[title={\ding{45}\quad Exercice~3.1. Classer des contrôles \niveaubase}]
Classez chacun des contrôles suivants (préventif, détectif ou correctif~;
manuel, automatisé ou semi-automatisé~; général ou d'application).
\begin{enumerate}[(a)]
    \item L'application refuse une facture dont le numéro de commande
          n'existe pas.
    \item Le responsable informatique revoit chaque trimestre la liste des
          comptes à privilèges.
    \item La comptable rapproche chaque mois le relevé bancaire et le grand
          livre.
    \item Une sauvegarde de la base est restaurée chaque mois sur un serveur
          de test.
    \item L'application produit chaque lundi la liste des ajustements de
          stock de la semaine, que la contrôleuse interne examine.
\end{enumerate}
\end{boxexo}

\begin{boxexo}[title={\ding{45}\quad Exercice~3.2. Les assertions \niveaubase}]
Pour chacun des risques suivants, indiquez l'assertion menacée et proposez
un contrôle.
\begin{enumerate}[(a)]
    \item Une livraison reçue le 2~janvier~2026 est enregistrée en
          décembre~2025.
    \item Une facture est saisie deux fois.
    \item Une vente est enregistrée sans livraison.
    \item Une réception n'est jamais saisie.
\end{enumerate}
\end{boxexo}

\begin{boxexo}[title={\ding{45}\quad Exercice~3.3. Le COSO et le SI \niveaumoyen}]
Pour les principes~3, 9, 11 et~16 du COSO, rédigez une question d'audit
portant sur le système d'information de la BANQUE SANAGA, et la preuve que
vous demanderiez pour y répondre.
\end{boxexo}

\begin{boxexo}[title={\ding{45}\quad Exercice~3.4. Corriger la matrice des droits \niveaumoyen}]
Reprenez la matrice des droits de LiZbiZ (section~\ref{sec:c03-sod}).
\begin{enumerate}[(a)]
    \item Proposez une nouvelle répartition des droits qui supprime toutes
          les paires incompatibles, sans retirer à quiconque les droits
          nécessaires à son métier.
    \item L'agence de Garoua ne compte que trois personnes. Proposez deux
          contrôles compensatoires pour le cumul F2 et F3 de M.~NANGA.
    \item Rédigez la ligne de la matrice des risques et des contrôles
          correspondant à l'un de ces contrôles.
\end{enumerate}
\end{boxexo}

\begin{boxexo}[title={\ding{45}\quad Exercice~3.5. La requête des conflits \niveaumoyen}]
La table \texttt{droits} est extraite de l'application, mais certains droits
sont accordés par un rôle~: une table \texttt{roles(utilisateur, role)} et
une table \texttt{role\_fonctions(role, fonction)}.
\begin{enumerate}[(a)]
    \item Écrivez la requête qui produit la liste complète des droits
          effectifs (directs et hérités).
    \item Modifiez la requête de la section~\ref{sec:c03-sod} pour qu'elle
          s'applique aux droits effectifs.
    \item Quelle vérification faites-vous avant d'exécuter ces requêtes~?
\end{enumerate}
\end{boxexo}

\begin{boxexo}[title={\ding{45}\quad Exercice~3.6. Tester un contrôle automatisé \niveaumoyen}]
L'application de LiZbiZ refuse un ajustement de stock supérieur à
500~000~FCFA sans validation.
\begin{enumerate}[(a)]
    \item Décrivez le test de conception et le test d'efficacité.
    \item Le journal des paramètres montre que le seuil a été porté à
          50~000~000~FCFA du 3 au 5~mars~2026, puis ramené à
          500~000~FCFA. Aucune demande de changement n'existe. Comment
          cotez-vous le contrôle sur la période, et pourquoi~?
\end{enumerate}
\end{boxexo}

\begin{boxexo}[title={\ding{45}\quad Exercice~3.7. \NON{} ou \NEK{}~? \niveaumoyen}]
Cotez chaque situation et justifiez.
\begin{enumerate}[(a)]
    \item Le critère exige une revue trimestrielle signée des droits~; le
          DAF affirme l'avoir faite, mais ne retrouve aucun document.
    \item Le journal des ajustements est remis, mais il s'arrête au
          15~mars~; la DSI explique qu'il a été « purgé ».
    \item Deux extractions de droits, faites à deux jours d'intervalle,
          donnent des listes différentes, sans changement enregistré entre
          les deux.
    \item Le contrôle n'existe dans aucune procédure, et personne ne
          l'exécute.
\end{enumerate}
\end{boxexo}

\begin{boxexo}[title={\ding{45}\quad Exercice~3.8. Lire un domaine \niveaumoyen}]
Seize contrôles des achats de LiZbiZ donnent 9~\OUI{}, 3~\NON{} et
4~\NEK{}. Calculez $\mu$, $U$ et $H$ (posez le calcul), puis comparez avec
les contrôles des stocks. Quelle démarche recommandez-vous dans chaque
cas~?
\end{boxexo}

\begin{boxexo}[title={\ding{45}\quad Exercice~3.9. Le dispositif de la BANQUE SANAGA \niveauexpert}]
La COBAC demande à la BANQUE SANAGA de compléter son dispositif de contrôle
interne sur la gestion des changements informatiques.
\begin{enumerate}[(a)]
    \item Identifiez trois risques liés aux changements du système central.
    \item Pour chacun, concevez un contrôle, en précisant sa nature, sa
          fréquence, son responsable et la preuve qu'il laisse.
    \item Pour l'un de ces contrôles, rédigez le test de conception et le
          test d'efficacité que réalisera l'audit interne.
    \item Indiquez dans quel cas l'auditrice devra coter le contrôle \NEK{}
          plutôt que \NON{}.
\end{enumerate}
\end{boxexo}

\begin{boxexo}[title={\ding{45}\quad Exercice~3.10. Contourner son propre contrôle \niveauexpert}]
Pour chacune des quatre lignes du Tableau~3.7, décrivez un moyen de
contourner le contrôle, puis le contrôle détectif qui révélerait ce
contournement. Concluez~: une matrice des risques et des contrôles peut-elle
être « complète »~?
\end{boxexo}

\bigskip

\begin{boxinfo}
\textbf{Pour aller plus loin.}
\begin{itemize}
    \item COSO, \textit{Internal Control -- Integrated Framework}
          (2013)~\cite{coso2013ic}~: le cadre lui-même, avec ses points
          d'attention par principe.
    \item COBAC, règlement R-2016/04~\cite{cobac2016r04}~: les exigences
          applicables aux banques de la CEMAC.
    \item Le guide d'audit des SI du CHAI~\cite{chai2014guide}~: des
          fiches de contrôles généraux directement utilisables.
\end{itemize}
\end{boxinfo}

\begin{boxinfo}
\textbf{Transition.} Un contrôle se justifie par le risque qu'il traite. Le
Chapitre~\ref{chap:risque} définit le risque, apprend à le mesurer et
présente le modèle du risque d'audit, point de départ de tout plan
d'audit.
\end{boxinfo}
